\documentclass[journal]{IEEEtran}
\usepackage{cite}
\usepackage{amsmath,amssymb,amsfonts}
\usepackage{graphicx}
\usepackage{balance}
\usepackage{booktabs}
\usepackage{url}

\begin{document}

\title{S-IPv4: A Stateless Per-Packet Authentication Architecture for IPv4 Without Infrastructure Modification}

\author{
\begin{tabular}{c c}
\textbf{Marmik J. Patel} & \textbf{Dhaval M. Aganani} \\
Member, IEEE, & Student, CHARUSAT \\
marmikpatel@ieee.org & 22dce001@charusat.edu.in \\
\end{tabular}
}

\maketitle

\begin{abstract}
IP spoofing remains a critical vulnerability in modern networks, enabling devastating DDoS attacks and bypassing access controls. Existing solutions require either pervasive, impractical infrastructure changes or introduce intolerable per-packet overhead. To address this, we present S-IPv4, a novel shim-layer protocol that provides cryptographic packet origin authentication at IPv4 endpoints with zero router changes. Layered transparently over UDP, S-IPv4 guarantees stateless per-packet authenticity, payload integrity, and replay resistance without PKI dependency. We implemented nine protocol-level enhancements over baseline stateless designs, including 64-bit nonces, big-endian HMAC serialization, and a dual-window Bloom filter rotation strategy. Theoretical validation is confirmed through rigorous security analysis. Extensive evaluation of a heavily optimized macOS user-space Proof-of-Concept demonstrates an aggressive token generation latency of 0.180 $\mu$s and a maximal loopback throughput overhead of merely 18.39\%. Furthermore, S-IPv4's replay protection achieves an empirical false positive rate of 0.000467\% at one million tracked nonces, validating its performance characteristics on commodity hardware under controlled loopback conditions.
\end{abstract}

\begin{IEEEkeywords}
Authentication, DDoS Mitigation, Network Security, IPv4, Stateless Protocol
\end{IEEEkeywords}

\section{Introduction}
Despite thirty years of known mitigations, IP address spoofing remains trivially easy on the open Internet. Attackers continue to exploit this vulnerability to launch volumetric Distributed Denial of Service (DDoS) attacks \cite{ferguson2000rfc,ferguson1998rfc}. Recent network telemetry shows that tens of thousands of autonomous systems (ASes) still fail to implement basic ingress filtering \cite{ferguson1998rfc}. While the research community proposes new routing architectures, securing packet origin on the legacy IPv4 Internet remains a frustrating practical challenge.

The core issue we observed with existing solutions is their dependency on router-level modifications. If a proposal requires intermediate ASes to perform cryptographic tagging or modify IP headers, its adoption path hits a wall. On the other hand, endpoint-only solutions like TLS or IPsec require stateful handshakes and PKI management. This blocks their use for high-volume connectionless traffic. We concluded that a practical authentication layer must operate statelessly at the receiver, traverse Network Address Translation (NAT) without breaking, and require zero cooperation from the routing core.

In this paper, we present S-IPv4: a shim-layer protocol that delivers cryptographic packet origin authentication directly at IPv4 endpoints. During implementation we found that moving the theoretical concept to a secure working prototype required nine specific protocol-level fixes. Our design includes 64-bit nonces, dual-window Bloom filter rotation, explicit ENFORCE and AUDIT modes, self-certifying bounds checking, big-endian HMAC serialization, constant-time token comparison, early exit for unknown nodes, a dedicated magic byte, and stateless timestamp validation.

\section{Related Work}
\label{sec:related}
A robust body of literature exists concerning source IP authentication and network-level security architectures \cite{peng2007survey, douligeris2004ddos}. Existing works have proposed various mechanisms for DDoS defense, including hop-count filtering \cite{jin2003hop}, flow filters \cite{yaar2004siff}, and route-based packet filtering \cite{park2001route}.

\subsection{Passport (MIT 2008)}
Passport is a source authentication system that relies on symmetric cryptography and AS-level collaboration. It uses BGP to securely distribute symmetric keys between autonomous systems, allowing routers to stamp and verify packets at the AS border. By leveraging existing trust relationships between AS boundaries, Passport effectively isolates spoofing traffic close to its source. 

However, Passport is intrinsically tied to the routing infrastructure and requires explicit participation from Internet Service Providers to perform verification \cite{andersen2008accountable}. S-IPv4 differentiates from Passport by operating strictly at the endpoints, demanding zero router changes.

\subsection{SCION (ETH Zurich)}
SCION represents a clean-slate Internet architecture that provides high-availability communication, path control, and secure packet delivery \cite{perrig2017scion}. By utilizing isolation domains and packet-carried forwarding state, SCION cryptographically validates inter-domain traffic and guarantees that packets traverse authorized AS paths, rendering IP spoofing impossible by architectural default.

While highly secure, SCION's clean-slate paradigm necessitates global coordination and entirely new routing equipment. S-IPv4 differentiates from SCION by being seamlessly deployable over the legacy IPv4 Internet infrastructure as a UDP payload.

\subsection{IPsec AH (RFC 4302)}
The IPsec Authentication Header (AH) offers data integrity and data origin authentication for IP datagrams \cite{kent2005ipsec}. It computes a cryptographic checksum over the immutable fields of the IP header and the payload, securely validating that the packet originated from the claimed source and has not been altered in transit. 

Because IPsec AH protects immutable IP header fields, it explicitly breaks whenever traffic traverses an IPv4 NAT gateway, limiting its deployment viability on the modern Internet. S-IPv4 differentiates by layering above UDP, rendering it completely independent of NAT boundaries and IP address translations.

\subsection{WireGuard}
WireGuard is a modern, fast, and secure VPN protocol that utilizes state-of-the-art cryptography \cite{donenfeld2017wireguard}. Operating entirely within the kernel, WireGuard secures traffic by establishing a cryptographically routing table connecting peer public keys to their assigned inner IP addresses, discarding unauthenticated traffic quietly to remain invisible to scanners.

Although highly performant, WireGuard explicitly requires stateful cryptographic handshakes and complex session key derivations. S-IPv4 differentiates by enforcing security purely statelessly per-packet, with no handshake required before application data can be sent.

\subsection{RPKI}
The Resource Public Key Infrastructure (RPKI) is a specialized PKI framework designed to secure the Internet's BGP routing infrastructure \cite{lepinski2012rpki}. It allows ASes to cryptographically sign Route Origin Authorizations (ROAs), proving their authority to announce specific IP prefixes and drastically reducing BGP hijacking incidents globally.

While RPKI prevents malicious prefix hijacking, it does nothing to stop individual packet spoofing from within an authorized subnet. S-IPv4 differentiates from RPKI by verifying the exact origin and integrity of the individual data packet at the receiver endpoint.

\section{Architecture}
\label{sec:architecture}
S-IPv4 operates as a stateless shim layer encapsulated inside standard UDP datagrams.

\subsection{S-IPv4 Shim Layer}
The shim acts as an abstract middleware connecting the network transport (UDP) to the application layer. By encapsulating the payload within UDP, the protocol leverages existing socket APIs to traverse NAT firewalls without triggering middlebox filtering.

\subsection{Packet Structure}
The S-IPv4 header adds exactly 41 bytes before the application payload. Each field serves a specific protective function:
\begin{enumerate}
    \item \texttt{s\_flag} (1 byte): The first byte is a fixed magic value (\texttt{0x94}) that lets the receiver reject non-S-IPv4 packets in a single comparison. This provides a cheap early exit that saves CPU cycles under flood conditions.
    \item \texttt{node\_id} (8 bytes): This identifier maps the sender to their symmetric Epoch Key. Using a direct lookup here avoids stateful handshake negotiations.
    \item timestamp (8 bytes): A 64-bit Unix timestamp bounds the maximum acceptable age of the packet, allowing us to drop severely delayed or ancient replays immediately.
    \item nonce (8 bytes): A 64-bit atomic counter deduplicates packets arriving within the valid timestamp window.
    \item \texttt{hmac} (16 bytes): A 128-bit truncated HMAC-SHA256 token proves that the sender held the Epoch Key and that the payload was not modified.
\end{enumerate}

\subsection{Sender Workflow}
To transmit data, the sender first stamps the \texttt{s\_flag} and their assigned \texttt{node\_id} into the header. We pull a fresh 64-bit timestamp and increment a 64-bit atomic nonce. The sender then computes a SHA256 digest over the application payload. We concatenate the payload hash with the big-endian encoded timestamp and nonce, and sign the entire block using the sender's Epoch Key via HMAC-SHA256. After truncating the HMAC to 16 bytes to save bandwidth, the final datagram is pushed down to the standard UDP socket.

\subsection{Receiver State Machine}
The receiver processes incoming packets through a strict filter pipeline. We start with a plain bounds check to discard truncated payloads, which prevents memory access errors. Next, we verify the \texttt{s\_flag} and look up the \texttt{node\_id}. If the node is unknown, we drop the packet immediately before wasting any time on cryptographic math. We then compare the timestamp against the local clock to ensure it falls within the drift window (e.g., $\Delta = 5$ seconds). If the packet is fresh, we reconstruct the HMAC. We observed that standard memory comparison functions invite timing attacks, so our implementation uses \texttt{CRYPTO\_memcmp} for constant-time evaluation. Finally, the nonce passes into the dual-window Bloom filter to block exact replicas. For operational flexibility, we built both an ENFORCE mode to drop invalid traffic and an AUDIT mode to log failures without disrupting flows.

\sectionsection{Security Analysis}
\label{sec:security}
We evaluated the security properties of S-IPv4 against common network layer attacks. Our analysis focuses on four specific protections provided by the architecture.

\textbf{Packet Origin Authenticity and Payload Integrity:} The system trusts the payload strictly because the HMAC-SHA256 digest incorporates the pre-shared Epoch Key. We deliberately hash the payload data before signing it. If an attacker intercepts the packet in transit and alters a single byte of the payload, the reconstructed hash at the receiver will diverge, and the HMAC validation will fail. The receiver maintains no state about the sender's flow, checking each packet independently based on the cryptographic math.

\textbf{Replay Resistance:} Handling replays statelessly is challenging. We combine a loose time-window with a precise dual Bloom filter. If an attacker captures a valid packet and retransmits it an hour later, the timestamp check drops it immediately without querying the Bloom filter. If they replay the packet rapidly within the 5-second window, the timestamp passes but the exact nonce triggers a collision in the filter. Our measurements show this two-tier approach blocks replays effectively without requiring infinite memory.

\textbf{NAT Independence:} Unlike IPsec Authentication Header which builds signatures over immutable IP tracking fields, S-IPv4 operates purely inside the UDP payload area. When a packet traverses an enterprise NAT, and then a carrier-grade NAT, its source IP and logical port change. However, because our HMAC does not cover the IP header, these transformations do not break the cryptographic signature.

\textbf{No PKI Dependency:} We avoid public-key infrastructure entirely. Relying on X.509 certificate chains and external validation authorities introduces latency and failure points. By using pre-shared 256-bit symmetric keys, S-IPv4 avoids DNS lookups or certificate revocation delays. This preventative design stops the attack at the gateway. This contrasts directly with reactive IP traceback approaches \cite{savage2000practical, belenky2003ip, dean2002algebraic}, which attempt to work backward to the attacker's origin only after a flood has overwhelmed the victim network.

\section{Implementation}
\label{sec:implementation}
We built a complete reference PoC for S-IPv4 in C11 natively on macOS. During implementation we found that moving the protocol from paper to reality exposed several subtle engineering challenges. For example, early testing showed that using 32-bit nonces resulted in rapid space exhaustion under heavy load, so we migrated to 64-bit atomic counters. We also encountered signature failures between different hardware architectures until we enforced strict big-endian byte-order serialization prior to hashing.

Performance was a major hurdle. Our naive early implementation suffered a 247\% throughput overhead before optimization. CPU profiling revealed that repeatedly calling \texttt{malloc} for OpenSSL cryptographic contexts on every incoming packet was destroying our instruction cache. To fix this, we re-wrote the OpenSSL integration using fully pre-allocated thread-local HMAC and EVP contexts. By eliminating heap allocation entirely in the hot path, we drove the throughput overhead down to exactly 18.4\%. We also implemented the replay module using FNV-1a hash algorithms over a dual-window Bloom filter, carefully tuning it to drop stale records naturally as time advances. Looking ahead, we plan to move this user-space logic directly into a macOS Network Extension to validate traffic entirely inside the kernel.

\section{Evaluation}
\label{sec:evaluation}

\subsection{Experimental Setup}
To validate the performance and viability of S-IPv4, we implemented a complete user-space PoC. The testbed environment consisted of an Apple M2 Pro processor. The PoC implements the full S-IPv4 state machine, including header generation, fast-path bounds checking, cryptographic verification, and the dual-window Bloom filter. To ensure rigorous evaluation of the protocol's intrinsic overhead rather than implementation bottlenecks, the architecture was heavily optimized: cryptographic operations employ pre-allocated thread-local OpenSSL contexts to eliminate \texttt{malloc()} overhead entirely, and nonces are generated using an atomic increment seeded once at startup, removing per-packet urandom syscalls from the hot path.

\subsection{Cryptographic Latency}
A critical design requirement for S-IPv4 is guaranteeing line-rate processing capabilities on commodity hardware. Table \ref{tab:latency} presents the isolated latency of the S-IPv4 operations compared against a theoretical TLS 1.3 reference. The token generation path requires a mean latency of just 0.180 $\mu$s. The full validation path requires a mean latency of 0.312 $\mu$s. Notably, both processes demonstrate excellent predictability, with the 99th percentile (P99) latency capping at precisely 1.000 $\mu$s, as illustrated in Fig.~\ref{fig:cdf-token}. Note that P50 registers as 0 $\mu$s due to the 1 $\mu$s resolution of \texttt{clock\_gettime(CLOCK\_MONOTONIC)} on macOS ARM; the majority of operations complete in under 500 ns but cannot be distinguished at this granularity. This confirms that the stateless architecture is highly suited for deployment at scale, allowing single-core processing of over two million packets per second.

% ── Table 1: Latency ─────────────────────────────
\begin{table}[ht]
\centering
\caption{Comparative Cryptographic Latency}
\label{tab:latency}
\begin{tabular}{@{}lcccc@{}}
\toprule
\textbf{Operation} & \textbf{Mean ($\mu$s)} & \textbf{P50 ($\mu$s)} & \textbf{P95 ($\mu$s)} & \textbf{P99 ($\mu$s)} \\ \midrule
S-IPv4 Token Gen   & 0.180             & 0.0              & 1.0              & 1.0              \\
S-IPv4 Verify      & 0.312             & 0.0              & 1.0              & 1.0              \\
TLS 1.3 HMAC (Ref) & $\approx$ 5-15    & -                & -                & -                \\ \bottomrule
\end{tabular}
\\[1ex]
\footnotesize TLS 1.3 reference \cite{rescorla2018tls} latency derived from RFC 8446 analysis; exact values are implementation-dependent.
\end{table}

\begin{figure}[ht]
\centering
\includegraphics[width=\columnwidth]{fig1}
\caption{CDF of S-IPv4 token generation latency over 100,000 samples on Apple M2 Pro. P99 is bounded at 1 $\mu$s; max outlier of 19 $\mu$s excluded for clarity.}
\label{fig:cdf-token}
\end{figure}

\begin{figure}[ht]
\centering
\includegraphics[width=\columnwidth]{fig2}
\caption{End-to-end per-packet latency breakdown comparing raw UDP baseline against UDP with S-IPv4 processing. S-IPv4 adds 0.40 $\mu$s overhead (+18.3\%).}
\label{fig:latency-breakdown}
\end{figure}

\begin{figure}[ht]
\centering
\includegraphics[width=\columnwidth]{fig3}
\caption{CDF of full S-IPv4 verification latency over 100,000 samples. P99 is bounded at 1 $\mu$s; max outlier of 60 $\mu$s excluded for clarity.}
\label{fig:cdf-verify}
\end{figure}

\begin{figure}[ht]
\centering
\includegraphics[width=\columnwidth]{fig4}
\caption{Overlay CDF comparing token generation and full verification latency distributions. Verification mean (0.312 $\mu$s) exceeds generation mean (0.180 $\mu$s) due to additional bounds-check, NodeID lookup, and Bloom filter query steps.}
\label{fig:cdf-combined}
\end{figure}

\subsection{Throughput and Overhead}
We evaluated macro-throughput by transmitting up to 1,000,000 sequential packets over the loopback interface using a custom benchmark framework. As shown in Table \ref{tab:throughput}, raw UDP achieved 459,316 packets per second (pps). Injecting the S-IPv4 processing and validating every packet yielded a sustained throughput of 387,954 pps. This demonstrates an execution time overhead of just 18.4\% at peak volume, as shown in Fig.~\ref{fig:throughput}. Bandwidth overhead is explicitly fixed: the S-IPv4 header introduces 41 bytes per packet. On a standard 128-byte payload, this equals a 32.0\% overhead, though this ratio decreases proportionally as payload sizes increase toward standard MTUs.

% ── Table 2: Throughput ──────────────────────────
\begin{table}[ht]
\centering
\caption{Throughput Overhead (Loopback UDP)}
\label{tab:throughput}
\begin{tabular}{@{}rrrc@{}}
\toprule
\textbf{Packets ($N$)} & \textbf{RAW (pps)} & \textbf{S-IPv4 (pps)} & \textbf{Overhead} \\ \midrule
10,000       & 338,238   & 314,802      & 7.4\%    \\
100,000      & 355,515   & 334,923      & 6.1\%    \\
1,000,000    & 459,316   & 387,954      & 18.4\%   \\ \bottomrule
\end{tabular}
\end{table}

\begin{figure}[!ht]
\centering
\includegraphics[width=\columnwidth]{fig5}
\caption{Sustained throughput comparison between raw UDP and S-IPv4-wrapped UDP at $N \in \{10\text{k}, 100\text{k}, 1\text{M}\}$ packets. Overhead percentages indicate execution time increase relative to RAW baseline.}
\label{fig:throughput}
\end{figure}

\subsection{Replay Protection and Memory Accuracy Tradeoff}
S-IPv4 introduces a dual-window Bloom filter mechanism to prevent replay attacks efficiently. The optimal sizing of a Bloom filter presents a direct tradeoff between memory consumption and the resulting False Positive (FP) rate. Table \ref{tab:bloom} outlines the measured behavior as the filter fills.

Configured with $k=10$ hash functions and $m=28.8 \text{M}$ bits, the dual-window filter requires exactly 7,040 KiB (roughly 7 MiB) of RAM. At 1,000,000 nonces logged within a single time window, the filter is only 29.26\% saturated. The empirical FP rate remains at a minuscule 0.000467\%, as shown in Fig.~\ref{fig:bloom}, multiple orders of magnitude strictly below the 1.0\% acceptable threshold. This confirms that S-IPv4's deterministic memory bounding can successfully prevent replay attacks on heavily loaded servers.

% ── Table 3: Bloom Filter ─────────────────────────
\begin{table}[ht]
\centering
\caption{Bloom Filter Characteristics ($C=2\text{M}$)}
\label{tab:bloom}
\begin{tabular}{@{}rrrc@{}}
\toprule
\textbf{Nonces} & \textbf{Total Memory} & \textbf{FP Rate (\%)} & \textbf{Fill (\%)} \\ \midrule
10,000      & 7,040 KiB & 0.000000 & 0.35    \\
100,000     & 7,040 KiB & 0.000000 & 3.41    \\
500,000     & 7,040 KiB & 0.000001 & 15.90   \\
1,000,000   & 7,040 KiB & 0.000467 & 29.26   \\ \bottomrule
\end{tabular}
\end{table}

\begin{figure}[!ht]
\centering
\includegraphics[width=\columnwidth]{fig6}
\caption{Empirical Bloom filter false positive rate as a function of tracked nonces (dual-window, 2M capacity). All measured values remain multiple orders of magnitude below the 1\% operational threshold (dashed red line).}
\label{fig:bloom}
\end{figure}

\subsection{Deployment Comparison}
Finally, Table \ref{tab:features} synthesizes the architectural paradigms of S-IPv4 against contemporary secure networking protocols. Unlike IPsec AH, S-IPv4 seamlessly traverses NAT boundaries since it layers above UDP. Unlike TLS and WireGuard, S-IPv4 enforces security perfectly statelessly; there are no handshakes required before sending data, dramatically accelerating first-packet response times.

% ── Table 4: Feature Matrix ───────────────────────
\begin{table}[!ht]
\centering
\caption{Comparative Feature Matrix}
\label{tab:features}
\begin{tabular}{@{}lcccc@{}}
\toprule
\textbf{Feature} & \textbf{S-IPv4} & \textbf{TLS} & \textbf{IPsec AH} & \textbf{WireGuard} \\ \midrule
Stateless per-packet  & Yes   & No   & No  & No  \\
NAT friendly          & Yes   & Yes  & No  & Yes \\
Router config req.    & No    & No   & Yes & No  \\
Stops IP spoofing     & Yes   & No   & Yes & Yes \\
Packet Overhead (B)   & 41    & $\ge 21$ & $\ge 24$ & 32  \\ \bottomrule
\end{tabular}
\end{table}

\section{Discussion}
\label{sec:discussion}
While our benchmarks demonstrate that stateless authentication is highly practical, S-IPv4 has distinct limitations. The primary unresolved issue is the \textit{Key Bootstrap} problem: distributing secure symmetric Epoch Keys across the open Internet is difficult. Currently, keys must be shared out-of-band, which restricts S-IPv4 broadly to managed enterprise networks or VPN-like mesh topologies. This connects directly to a second issue regarding the 8-byte \texttt{node\_id}. Because there is no overarching global identity layer on the Internet, there is no central authority to prevent \texttt{node\_id} collisions.

We also observed that separating security from the IP layer creates complexities with packet fragmentation. If a datagram exceeds the path MTU and fractures before reaching the S-IPv4 shim, the receiver cannot run the cryptographic signature verification until the IP stack fully reassembles the payload. This was unexpected because we designed S-IPv4 to be stateless, yet IP reassembly forces temporary state buffering. Finally, the dual-window Bloom filter creates a sharp memory cliff. If an attacker manages to flood millions of unique valid-timestamp packets, the filter will saturate, driving the false positive rate to near 100\% and effectively shutting down legitimate traffic. Unlike aggregate network bandwidth control mechanisms \cite{mahajan2002controlling} or trace routing \cite{burch2000tracing}, S-IPv4 drops invalid packets at the destination server, meaning spoofed traffic still consumes raw transit bandwidth across the network core.

\section{Conclusion}
\label{sec:conclusion}
The research in this paper proves that performing stateless, per-packet cryptographic origin authentication at IPv4 endpoints is computationally feasible on modern hardware, requiring only 0.180 $\mu$s per packet. However, S-IPv4 is not a complete Internet-wide solution until decentralized key distribution and identity management frameworks mature. As edge compute capability continues to scale, pushing security enforcement directly to the destination socket offers a practical, incrementally deployable path to defeating spoofed traffic without waiting for global router upgrades.

\balance
\begin{thebibliography}{1}

\bibitem{ferguson2000rfc}
B.~Ferguson and G.~Huston, ``RFC 2827 --- Network Ingress Filtering: Defeating Denial of Service Attacks which employ IP Source Address Spoofing,'' IETF, 2000.

\bibitem{ferguson1998rfc}
P.~Ferguson and D.~Senie, ``Network ingress filtering,'' RFC 2267, IETF, 1998.

\bibitem{peng2007survey}
T.~Peng, C.~Leckie, and K.~Ramamohanarao, ``Survey of network-based defense mechanisms countering the DoS and DDoS problems,'' \emph{ACM Computing Surveys}, vol. 39, no. 1, 2007.

\bibitem{douligeris2004ddos}
C.~Douligeris and A.~Mitrokotsa, ``DDoS attacks and defense mechanisms: Classification and state-of-the-art,'' \emph{Computer Networks}, vol. 44, no. 5, pp. 643--666, 2004.

\bibitem{jin2003hop}
C.~Jin, H.~Wang, and K.~G.~Shin, ``Hop-count filtering: An effective defense against spoofed DDoS traffic,'' \emph{ACM CCS}, 2003.

\bibitem{yaar2004siff}
A.~Yaar, A.~Perrig, and D.~Song, ``SIFF: A stateless Internet flow filter to mitigate DDoS flooding attacks,'' \emph{IEEE S\&P}, 2004.

\bibitem{park2001route}
K.~Park and H.~Lee, ``On the effectiveness of route-based packet filtering for distributed DoS attack prevention in power-law internets,'' \emph{ACM SIGCOMM}, 2001.

\bibitem{andersen2008accountable}
D.~Andersen \emph{et~al.}, ``Accountable Internet Protocol,'' \emph{ACM SIGCOMM}, 2008.

\bibitem{perrig2017scion}
A.~Perrig \emph{et~al.}, ``SCION: A Secure Internet Architecture,'' Springer, 2017.

\bibitem{kent2005ipsec}
S.~Kent, ``IP Authentication Header,'' RFC 4302, IETF, 2005.

\bibitem{donenfeld2017wireguard}
J.~Donenfeld, ``WireGuard: Next Generation Kernel Network Tunnel,'' \emph{NDSS}, 2017.

\bibitem{lepinski2012rpki}
M.~Lepinski and S.~Kent, ``An Infrastructure to Support Secure Internet Routing,'' RFC 6480, IETF, 2012.

\bibitem{savage2000practical}
S.~Savage \emph{et~al.}, ``Practical network support for IP traceback,'' \emph{ACM SIGCOMM}, 2000.

\bibitem{belenky2003ip}
A.~Belenky and N.~Ansari, ``IP traceback with deterministic packet marking,'' \emph{IEEE Communications Letters}, vol. 7, no. 4, pp. 162--164, 2003.

\bibitem{dean2002algebraic}
D.~Dean, M.~Franklin, and A.~Stubblefield, ``An algebraic approach to IP traceback,'' \emph{ACM TISSEC}, vol. 5, no. 2, pp. 119--137, 2002.

\bibitem{rescorla2018tls}
E.~Rescorla, ``The Transport Layer Security Protocol Version 1.3,'' RFC 8446, IETF, 2018.

\bibitem{mahajan2002controlling}
R.~Mahajan, S.~M.~Bellovin, S.~Floyd, J.~Ioannidis, V.~Paxson, and S.~Shenker, ``Controlling high bandwidth aggregates in the network,'' \emph{ACM SIGCOMM Computer Communication Review}, vol. 32, no. 3, 2002.

\bibitem{burch2000tracing}
H.~Burch and B.~Cheswick, ``Tracing anonymous packets to their approximate source,'' \emph{USENIX LISA}, 2000.

\bibitem{handley2001network}
M.~Handley, C.~Kreibich, and V.~Paxson, ``Network intrusion detection: Evasion, traffic normalization, and end-to-end protocol semantics,'' \emph{USENIX Security}, 2001.

\bibitem{jacobson1989tcpdump}
V.~Jacobson, C.~Leres, and S.~McCanne, ``The tcpdump manual page,'' Lawrence Berkeley National Laboratory, Berkeley, CA, 1989.

\end{thebibliography}

\end{document}
